% Options for packages loaded elsewhere
\PassOptionsToPackage{unicode}{hyperref}
\PassOptionsToPackage{hyphens}{url}
\documentclass[
  11pt,
]{article}
\usepackage{xcolor}
\usepackage[margin=1in]{geometry}
\usepackage{amsmath,amssymb}
\setcounter{secnumdepth}{-\maxdimen} % remove section numbering
\usepackage{iftex}
\ifPDFTeX
  \usepackage[T1]{fontenc}
  \usepackage[utf8]{inputenc}
  \usepackage{textcomp} % provide euro and other symbols
\else % if luatex or xetex
  \usepackage{unicode-math} % this also loads fontspec
  \defaultfontfeatures{Scale=MatchLowercase}
  \defaultfontfeatures[\rmfamily]{Ligatures=TeX,Scale=1}
\fi
\usepackage[]{libertinus}
\ifPDFTeX\else
  % xetex/luatex font selection
\fi
% Use upquote if available, for straight quotes in verbatim environments
\IfFileExists{upquote.sty}{\usepackage{upquote}}{}
\IfFileExists{microtype.sty}{% use microtype if available
  \usepackage[]{microtype}
  \UseMicrotypeSet[protrusion]{basicmath} % disable protrusion for tt fonts
}{}
\usepackage{setspace}
\makeatletter
\@ifundefined{KOMAClassName}{% if non-KOMA class
  \IfFileExists{parskip.sty}{%
    \usepackage{parskip}
  }{% else
    \setlength{\parindent}{0pt}
    \setlength{\parskip}{6pt plus 2pt minus 1pt}}
}{% if KOMA class
  \KOMAoptions{parskip=half}}
\makeatother
\usepackage{longtable,booktabs,array}
\usepackage{caption}
\captionsetup[table]{skip=6pt}
\newcounter{none} % for unnumbered tables
\usepackage{calc} % for calculating minipage widths
% Correct order of tables after \paragraph or \subparagraph
\usepackage{etoolbox}
\makeatletter
\patchcmd\longtable{\par}{\if@noskipsec\mbox{}\fi\par}{}{}
\makeatother
% Allow footnotes in longtable head/foot
\IfFileExists{footnotehyper.sty}{\usepackage{footnotehyper}}{\usepackage{footnote}}
\makesavenoteenv{longtable}
\setlength{\emergencystretch}{3em} % prevent overfull lines
\providecommand{\tightlist}{%
  \setlength{\itemsep}{0pt}\setlength{\parskip}{0pt}}
\usepackage{bookmark}
\IfFileExists{xurl.sty}{\usepackage{xurl}}{} % add URL line breaks if available
\urlstyle{same}
% Make links footnotes instead of hotlinks:
\DeclareRobustCommand{\href}[2]{#2\footnote{\url{#1}}}
% fallback for those not using the hyperref driver hyperxmp:
\makeatletter
\@ifundefined{xmpquote}{\newcommand{\xmpquote}[1]{#1}}{}
\makeatother
\hypersetup{
  hidelinks,
  pdfcreator={LaTeX via pandoc}}

\author{}
\date{}

\begin{document}

\setstretch{1.2}
\section{Abstract}\label{abstract}

Why did the fossil economy begin in Britain rather than in the
Netherlands, with Europe's densest waterways, or China, with coal and
rivers in abundance? Terje Tvedt's answer, in this journal, was that
Britain alone could engineer its water to deliver coal cheaply: water
was the precondition for a coal economy. That is a claim about sequence,
and Britain is where the annual series to test it exist together. Using
British sectoral output for 1700--1870, a new canal-mileage series,
installed horsepower and the Google Books corpus, we find two growth
regimes: an aggregate acceleration from 1775--1792, as the canals were
built, absorbed by population, and a per-capita acceleration from 1818
that belongs to steam. Canal mileage predicts coal output, not income.
At the Maddison benchmarks the Dutch, Belgian, Chinese and Indian
records are consistent with the rule that neither engineered water nor
coal was sufficient alone.

\textbf{Keywords:} Industrial Revolution, canals, coal, water
infrastructure, Great Divergence, Netherlands, China, energy transition,
growth regimes

\begin{center}\rule{0.5\linewidth}{0.5pt}\end{center}

\section{1. Introduction}\label{introduction}

``Why England and not China and India?'' Terje Tvedt asked in this
journal in 2010, and answered with water (Tvedt 2010). Pomeranz had
asked the same question of the Yangzi delta a decade earlier and
answered with coal and colonies (Pomeranz 2000). The two answers are
usually read as rivals. They need not be. In 1700 the parts of what
became Britain's fossil economy were scattered across several economies,
and none held them all. The Netherlands had the densest network of
engineered waterways in Europe and no coal (de Vries 1978; de Zeeuw
1978). China had coal in quantity and the longest artificial waterway in
the world, and the two lay a thousand kilometres apart (Pomeranz 2000;
Elvin 1973). India had coal in Bengal and its cotton industry on the
western and south-eastern coasts, joined by rivers whose flow followed
the monsoon (Parthasarathi 2011; Tvedt 2010). Belgium had coal under
Hainaut and Liège and would acquire the canals to move it in 1818 and
1832, after which its income per head grew faster than Britain's (Mokyr
1976). Britain had coal, rivers with the gradient to drive mills and the
flow to float barges, and a legal form, the canal Act, that let private
capital join the two. By 1830 no economy had gone further in running on
mineral fuel moved by water.

Tvedt's claim is that the water system was the \emph{precondition} for
the coal economy, not a rival to it. This paper takes that claim
literally and tests it. A precondition is a claim about order and
dependence, and it makes predictions that a rival-cause claim does not.
It predicts that the sectors served by canals accelerated before steam
was a significant source of power; that coal output responded to the
canal network; that steam capacity followed coal; and that the sectors
canals did not serve, and the outcome canals could not directly raise,
did not move. Each prediction is about sequence within one economy, and
sequence can be tested only where the relevant annual series exist
together. The Netherlands and Sweden have annual output from before
1700, but only Britain also has canal mileage and installed power year
by year over a period in which all three changed. The comparative cases
can then be read through its result.

The test has to reconcile two clocks that the literature has kept apart.
The steam clock starts in the 1780s with Watt's rotative engine and
reads industrialisation as the arrival of a new prime mover and of
growth in income per head, which barely moved before 1830 (Landes 1969;
Allen 2009; Crafts 1985; Crafts and Harley 1992). The canal clock starts
in 1761, when the Duke of Bridgewater's canal halved the price of coal
in Manchester, and is calibrated to tonnage, prices and towns (Bogart
2014; Szostak 1991; Maw 2013). Setting the two side by side on annual
data produces three findings.

First, Britain had two growth regimes, not one. Using the annual output
series of Broadberry, Campbell, Klein, Overton and van Leeuwen, we find
trend breaks between 1775 and 1792 in total output, industry, services,
coal, iron and population, and a break in 1818 in income per head. In
the first regime the growth of aggregate output nearly tripled and the
growth of population quadrupled, while income per head grew no faster
than it had since 1700 and, between 1760 and 1790, slower. Population
absorbed the acceleration. Both clocks are right, about different
variables.

Second, the canal network is a dose that predicts the right things. We
assemble a year-by-year series of canal miles from the completion dates
of 155 English and Welsh canals and corroborate its timing with 152
parliamentary authorisations recovered from Priestley's 1831 survey.
Cumulative canal mileage predicts coal output over the following five to
twenty years and it predicts population; it predicts neither income per
head nor agricultural output, and the coal result survives a quadratic
trend, first differences and two predetermined versions of the dose.
Steam horsepower, in turn, raises income per head only after 1830. In
the Google Books British corpus the same order appears, to within a
decade: ``coal barge'' and ``coal wharf'' reach a quarter of their
mid-century frequency around 1781 and 1800; ``steam engine'' and ``steam
power'' around 1808 and 1826.

Third, the comparative cases behave as a precondition thesis says they
should, and the standard cross-country test cannot see it. Between the
Maddison benchmark years 1700 and 1820, before steam was a third of
British power, British output tripled while Dutch output grew by 9 per
cent: water transport without coal, and without falling water to drive
mills, did not compound. Belgium, whose coalfield canals opened in 1818
and 1832, grew per head by 8 per cent between 1700 and 1820 and by 82
per cent in the next fifty years, more than Britain. China nearly
trebled its population and lost income per head. A cross-country
difference-in-differences on income per head, meanwhile, locates a
British take-off in 1807, which is the year the Dutch economy lost 44
per cent of its income under French occupation while Britain held level.

The contribution is to four literatures. To global history and the Great
Divergence debate, we give Tvedt's thesis a quantitative test and
restate it as a rule the comparative cases can be checked against:
neither engineered water nor coal was sufficient, and the economies that
had one without the other did not industrialise until they had both. To
the transport-revolution literature, we add an annual national dose
series and a sequencing test that connects canals to the energy
transition rather than to urban growth, the channel recently quantified
by Alvarez-Palau and co-authors (Alvarez-Palau et al.~2025). To the
energy-history debate among Wrigley, Allen, Pomeranz and Malm, we add
the observation that the escape from the Malthusian ceiling was under
way, in aggregate, before steam, in Wrigley's advanced organic economy.
To the practice of historical difference-in-differences (Roth et
al.~2023), we add a worked example of how a control group's catastrophe
becomes a treatment group's miracle.

The next section places the argument in the literature; the following
two describe the data and report the results, British first and
comparative second; the discussion returns to Tvedt's question, and the
limitations, which are real, come before the conclusion.

\begin{center}\rule{0.5\linewidth}{0.5pt}\end{center}

\section{2. Water, coal and the divergence
debate}\label{water-coal-and-the-divergence-debate}

\subsection{2.1 The transport revolution and its
canals}\label{the-transport-revolution-and-its-canals}

The transport revolution of eighteenth-century Britain is well
documented and, until recently, poorly measured. Between the Bridgewater
Canal of 1761 and the opening of the Liverpool and Manchester Railway in
1830, Parliament authorised roughly 165 canal companies and the
navigable network grew from about 1,400 miles of improved river to some
4,000 miles of river and canal (Hadfield 1984; Bogart 2014). The coal
trade itself was older and larger than the canals: the coastal trade
from the Tyne and Wear to London was already the largest bulk-freight
flow in Europe in 1700, and Britain's coal output had been rising for
two centuries before the first canal was cut (Hatcher 1993; Flinn 1984;
Pollard 1980). What canals changed was inland access. Priestley's survey
of 1831 opens its account of the Bridgewater Canal by noting that the
Duke's ``primary object'' was ``to open his valuable collieries at
Worsley, and to supply the town of Manchester with coal, at a much
cheaper rate'' (Priestley 1831). Turnbull's regional study made the same
point: canals were built where coal was, carried coal above all else,
and the counties that acquired them saw coal and population grow faster
(Turnbull 1987).

Quantitative work has caught up in the last decade. Bogart's survey
collects the mileage, cost and freight-rate evidence and concludes that
inland freight costs fell by half or more between 1700 and 1830 (Bogart
2014). Alvarez-Palau, Bogart, Satchell and Shaw-Taylor estimate that
transport costs relative to producer prices fell by nearly 75 per cent
between 1680 and 1830 and that, without that fall, inland towns would
have been 20 to 25 per cent smaller in 1841 (Alvarez-Palau et al.~2025).
Their outcome is urban population. Ours is the energy system. The two
are complementary: towns grew where coal could be delivered, and coal
could be delivered where water had been engineered.

\subsection{2.2 Energy, the organic economy and the Great
Divergence}\label{energy-the-organic-economy-and-the-great-divergence}

The energy literature starts from a different place. Wrigley's organic
economy is bounded by the annual product of the land; sustained growth
required a ``mineral-based energy economy'' in which coal replaced wood
as a source of heat and, much later, water and wind as a source of power
(Wrigley 1988; Wrigley 2010; Wrigley 2016). Wrigley is careful that the
first transition was under way from the sixteenth century and that the
eighteenth-century economy was an advanced organic one, burning mineral
heat while its power remained organic; Warde's energy accounts confirm
the chronology (Warde 2007). The steam clock belongs to Wrigley's
readers more than to Wrigley. Allen's induced-innovation account makes
cheap British coal and dear British labour the reason the steam engine
was worth inventing in Britain and nowhere else (Allen 2009). Pomeranz
gives coal a starring role in the Great Divergence, alongside the ghost
acres of the Americas (Pomeranz 2000). Clark and Jacks and Fernihough
and O'Rourke have found that proximity to coal mattered a great deal for
where Europe industrialised (Clark and Jacks 2007; Fernihough and
O'Rourke 2021). Malm's \emph{Fossil Capital} accepts coal's centrality
but denies that its adoption was driven by thermodynamic superiority or
scarcity of water power: steam won because it freed capital from the
riverbank and allowed mills to be sited where labour could be
disciplined (Malm 2016).

What these accounts share is a clock calibrated to steam and coal
combustion. Tvedt's intervention is to move the clock upstream. In ``Why
England and not China and India?'', he argues that Britain's distinctive
endowment was not coal as such, which China had in abundance, but a
hydrological regime of gentle gradients, steady flow and navigable
rivers that could be engineered into a national network of canals and
water-powered mills at low cost (Tvedt 2010). The monsoonal rivers of
Asia could not. On this reading the water system is the precondition: it
created the market integration, the cheap coal at the point of use, and
the water-powered factory template on which steam later built. The Asian
side of the comparison has its own literature. Elvin's high-level
equilibrium trap describes a Chinese economy that had exhausted its
organic frontier without reaching its mineral one, and Pomeranz locates
the reason in geography: China's coal lay in the north-west, far from
the Jiangnan core (Elvin 1973; Pomeranz 2000). Parthasarathi makes a
parallel case for India, whose coal in Bengal lay far from the cotton
regions and whose rivers ran with the monsoon, and argues that
endowments and pressures, not institutions or culture, pushed Britain to
coal (Parthasarathi 2011). Tvedt's evidence is comparative and
qualitative. This paper asks whether the annual British record supports
the sequence he proposes, and then reads the comparative cases against
what it finds.

A precondition is not a rival cause. Coal remains the fuel of the second
regime and steam its engine. The claim is that neither would have
arrived when and where it did without the first regime, and that the
first regime is visible in the data if one looks at the variables canals
could move: coal, aggregate output, and the number of people the economy
could feed.

\subsection{2.3 Two chronologies of
growth}\label{two-chronologies-of-growth}

The macroeconomic chronology has moved decisively towards gradualism.
Crafts, and Crafts and Harley, cut the growth rates of the 1780--1830
period that Deane and Cole had estimated and showed that income per head
grew slowly until well into the nineteenth century (Crafts 1985; Crafts
and Harley 1992). Broadberry, Campbell, Klein, Overton and van Leeuwen's
reconstruction of British output from the medieval period confirms the
picture: annual growth of GDP per head averaged around 0.3 per cent from
1700 to the 1820s and only then accelerated (Broadberry et al.~2015).
Crafts's growth-accounting of steam finds its contribution to
labour-productivity growth trivial before 1830 and dominant only after
1850 (Crafts 2004). On Kanefsky's estimates of installed power, steam
did not overtake water until the 1830s (Kanefsky 1979; Kanefsky and
Robey 1980).

Gradualism in income per head has been read as a verdict against any
eighteenth-century take-off. It is not. The same Broadberry series show
total output and population doubling between 1760 and 1830. A growth
regime in which aggregate output accelerates while income per head does
not is exactly what Malthus described and exactly what an infrastructure
shock in an organic economy should produce: more mouths fed, not richer
mouths. The distinction between the aggregate and the per-capita clock
is, we argue, the key to reconciling the canal historians with the
cliometricians.

\subsection{2.4 Text as data and historical
difference-in-differences}\label{text-as-data-and-historical-difference-in-differences}

Two methodological literatures frame our tools. Michel et al.~introduced
the Google Books corpus as a quantitative record of culture, and
Pechenick, Danforth and Dodds documented its limits: it over-weights
technical publishing and records what was printed rather than said
(Michel et al.~2011; Pechenick et al.~2015). For infrastructure
vocabulary that bias is a feature. The historical
difference-in-differences literature has grown cautious about serial
correlation, heterogeneous effects and pre-trends (Bertrand, Duflo, and
Mullainathan 2004; Roth et al.~2023; Rambachan and Roth 2023); we add a
more elementary hazard, a control group hit by a war the treatment group
escaped, whose economic consequences Crouzet described sixty years ago
(Crouzet 1964).

\begin{center}\rule{0.5\linewidth}{0.5pt}\end{center}

\section{3. Data and empirical
strategy}\label{data-and-empirical-strategy}

\subsection{3.1 British output, population and
power}\label{british-output-population-and-power}

Annual real output for Great Britain by sector, 1700--1870, comes from
Broadberry, Campbell, Klein, Overton and van Leeuwen as published in the
Bank of England's \emph{A Millennium of Macroeconomic Data for the UK},
version 3.1 (Broadberry et al.~2015; Thomas and Dimsdale 2017). We use
total GDP, agriculture, industry and services, and from the
industrial-production tables the physical output indices for coal, iron
and textiles. Population of Great Britain is from the same source, based
on Wrigley and Schofield's reconstitution to 1801 and the censuses
thereafter. GDP per head is total output divided by population. The
series are indices with 1700 = 100.

Installed stationary power is from Kanefsky's estimates as tabulated by
Crafts, at the benchmark years 1760, 1800, 1830 and 1870, reproduced in
Table 4 (Kanefsky 1979; Crafts 2004). We interpolate log-linearly
between the four benchmarks, wind included. The interpolated series are
smooth by construction, and we treat inference that depends on their
year-to-year variation with corresponding caution.

\subsection{3.2 The canal network as a
dose}\label{the-canal-network-as-a-dose}

A single treatment date cannot represent the canal network. It was built
over sixty years in two waves, and a binary indicator cannot distinguish
the Grand Cross of the 1770s from the mania canals completed between
1790 and 1816. We therefore construct a continuous dose, which we call
the canal stock: cumulative canal miles in operation by year.

The primary series is assembled from the completion year and length of
155 canals in England and Wales, taken from the reference table of
British canals maintained on Wikipedia and retrieved on 3 September
2026; we include every entry with a stated length and opening year, and
date staged openings to full completion, which biases the series late by
a few years. The 155 canals total 2,967 miles, of which 2,320 were open
by 1830. The series excludes river navigations improved before 1700,
which were the bulk of the network's initial 1,400 miles, so it measures
additions to the network rather than its level; for the growth
regressions this is the relevant quantity. The Cambridge Group's dynamic
GIS puts the England and Wales network at about 1,400 miles in 1760 and
4,000 in 1835, most of the growth canal (Satchell 2017); our 2,102 miles
over 1760--1830 are four-fifths of that increase, the rest post-1830
completions, branches and river work we omit. Hadfield's roughly 165
canal companies authorised between 1761 and 1830 counts undertakings,
some of which built no canal, and is not comparable (Hadfield 1984).

We corroborate the timing with an independent source. Priestley's
\emph{Historical Account of the Navigable Rivers, Canals, and Railways
throughout Great Britain} records the Acts that authorised each
navigation with regnal year and date of Royal Assent (Priestley 1831).
From the digitised 1831 text we identify 325 entries by their
capitalised headings and recover a first-authorisation year for 152,
including 70 canals authorised between 1760 and 1829, thirty of them in
the 1790s; the rest are lost to the quality of the optical character
recognition, and lengths could not be recovered at all. The
authorisation series leads the completion series by roughly a decade, as
it should, and both show the same two waves.

\subsection{3.3 Comparative benchmarks}\label{comparative-benchmarks}

Cross-country GDP per head and population are from the Maddison Project
Database 2023 (Bolt and van Zanden 2025). The series labelled GBR is the
United Kingdom including Ireland; we note where this matters. For
Britain, France, the Netherlands, Sweden, Germany, Spain and Portugal
the series is annual from 1700; for Belgium, China and India it is
interpolated between benchmarks. Population in the Maddison database is
a benchmark series before 1820, so the comparative cases are read at the
benchmark years 1700, 1820 and 1870, alongside what the national
literatures record about fuel, gradient and waterways: de Vries and de
Zeeuw for the Netherlands, Mokyr for Belgium, Pomeranz and Elvin for
China and Parthasarathi for India (de Vries 1978; de Zeeuw 1978; Mokyr
1976; Pomeranz 2000; Elvin 1973; Parthasarathi 2011).

\subsection{3.4 Text data}\label{text-data}

Word and phrase frequencies are from the Google Books Ngram corpus,
British English 2019 edition, 1700--1900, retrieved with the viewer's
three-year smoothing (Michel et al.~2011). The unigram ``canal'' is used
to test what the print record measures. For the semantic sequence we
retrieved sixteen bigrams naming coal's mode of use and use eight:
``coal barge'', ``coal wharf'', ``coal boat'' and ``canal boat'' for
coal moved by water; ``steam engine'' and ``steam power'' for coal
burned in engines; ``coal field'' as the geologist's term; and ``fire
engine'', the eighteenth-century name for an atmospheric engine, as a
check on terminological drift. ``Steam boat'' was retrieved and excluded
because it names water transport as much as steam. Each bigram is
indexed to its own 1850 frequency and the groups are equal-weighted
averages of their members, so that a single common term cannot dominate;
a further five-year centred mean is applied before thresholds are read.

\subsection{3.5 Empirical strategy}\label{empirical-strategy}

Our questions are about timing, sequence and incidence within one
national economy over 170 years. The sample is small, the series are
smooth and trending, and the treatment is itself a slowly accumulating
stock. We match the tools to those facts, and we use linear methods
throughout: with 130 annual observations, flexible machine-learning
estimators fit the time trend closely enough to leave no treatment
variation to estimate from.

\textbf{Regime detection.} For each British series we estimate a linear
trend in logs and test for a single break in level and slope at an
unknown date using the Andrews sup-F statistic with 15 per cent
trimming, reporting the statistic against the Andrews critical value,
and for two breaks by grid search with a minimum segment of twenty years
(Andrews 1993; Bai and Perron 1998). We also fix the break at 1761 and
test for a change in trend slope over 1700--1830 with Newey--West
standard errors (Newey and West 1987); since the data-chosen breaks fall
later, this fixed date is a conservative pre-test. Agriculture and GDP
per head are the placebo series: canals did not carry harvests to any
great extent, and a precondition operating through population should
leave income per head unchanged.

\textbf{Dose-response.} We regress the log of each output series on the
canal stock, in thousands of miles, and a linear trend over 1700--1830,
the period before railways. Because a cumulative stock is itself smooth
and accelerating, we report three further specifications: a quadratic
trend with a dummy for war years (1756--63, 1775--83, 1793--1815); a
first-difference model with ten lags of new mileage, reporting the sum
of the lag coefficients and the joint F-test; and a horse race that adds
the print frequency of ``steam'' and the war dummy. We test the
residuals of the level regressions for a unit root. We build mainly on
coefficients that survive all specifications; where an outcome survives
most but not all, we say so.

\textbf{Predetermined doses.} Canals were promoted where trade was
growing, so the contemporaneous stock is not exogenous. We therefore
also use two doses fixed before the outcome window: the completion-based
stock lagged ten and fifteen years, and the cumulative number of canals
\emph{authorised} by Act of Parliament, from Priestley, lagged ten
years. An Act precedes completion by roughly a decade and cannot respond
to output after it is passed. These doses remove contemporaneous
feedback; they do not remove anticipation, and we test for it directly
by regressing new authorisations on past output growth and by adding the
outcome's own past growth to the dose regressions.

\textbf{Sequencing by local projection.} To test the chain water → coal
→ steam → income per head, we estimate local projections (Jordà 2005):
the change in the log outcome over horizons of one to twenty years is
regressed on the regressor of interest, the outcome's own level and a
trend, with Newey--West errors whose bandwidth is the larger of ten
years and the horizon. Agriculture is again the placebo. Steam equations
start in 1760, when the horsepower series begins; the income equations
are estimated on 1760--1830, 1760--1870 and 1830--1870, with an
interaction of horsepower and a post-1830 indicator, since the thesis
predicts that steam's per-capita effect appears only after it becomes
the majority power source. With eight outcomes, four specifications and
four reported horizons, we treat a Bonferroni threshold of about 0.002
as the family-wise bar and say which results clear it.

\textbf{Semantic sequencing.} For each bigram group we record the first
year in which the five-year moving average reaches 10, 25 and 50 per
cent of its 1850 level, and the year in which ``steam engine'' overtakes
``fire engine'' in print. Because the British sub-corpus is thin before
1780, we repeat the 25 per cent threshold with three-, five- and
nine-year windows and with 1830 as well as 1850 as the reference year.
We also re-estimate the steam links of the chain with the print
frequency of ``steam engine'' in place of interpolated horsepower, as an
independent, if noisier, measure of steam's diffusion.

\textbf{Comparative cases and the cross-country test.} The comparative
cases are read as four checks on the rule the British result supplies:
an economy with engineered water and no coal, one with coal and no
engineered water, one that acquired both late, and one with neither in
reach. We report growth between the Maddison benchmark years for each.
To show why we do not go further, we also estimate the two-way
fixed-effects difference-in-differences on Maddison GDP per head, in
levels and logs, with 1761 as treatment date, for three control groups
and three sample windows, with its event-study version and clustered and
collapsed standard errors, and compare peak-to-trough drawdowns over
1785--1815 by country.

All code, the tidy data files, the Priestley parsing rules and the
horsepower interpolation are in the replication repository named in the
data availability statement, together with a mediation decomposition
that is not identified once a quadratic trend is absorbed and is
reported there as a negative result. The code and the text were prepared
with the assistance of an AI tool under the author's direction, as
declared in the acknowledgements.

\begin{center}\rule{0.5\linewidth}{0.5pt}\end{center}

\section{4. Results}\label{results}

\subsection{4.1 Two growth regimes}\label{two-growth-regimes}

Figure 1 plots the British series on a logarithmic scale with the two
canal-building waves shaded. Panel (a) shows total output, population
and income per head; panel (b) shows coal, industry and agriculture. The
visual impression is of a fan opening after 1760: total output and
population steepen together while income per head continues at its
previous slope until the 1820s. Agriculture never steepens at all.

Figure 1: Britain's two growth regimes. Annual indices, 1700 = 100, log
scale. Shaded bands mark the first canal wave (1760--1780) and the
canal-mania completions (1790--1816). Vertical lines at 1761 and 1818.
Source: Bank of England millennium dataset (Broadberry et al.~2015).

Table 1 gives the trend growth rates by period. Between 1700--1760 and
1790--1815, the growth of total output nearly tripled, coal's more than
tripled and population's quadrupled. Income per head grew at 0.25 per
cent a year in the first period, 0.08 in 1760--1790 and 0.40 in
1790--1815. Agriculture is flat throughout.

\textbf{Table 1: Trend growth of British series, per cent per year}

{\def\LTcaptype{none} % do not increment counter
\begin{longtable}[]{@{}
  >{\raggedright\arraybackslash}p{(\linewidth - 14\tabcolsep) * \real{0.1250}}
  >{\raggedleft\arraybackslash}p{(\linewidth - 14\tabcolsep) * \real{0.1250}}
  >{\raggedleft\arraybackslash}p{(\linewidth - 14\tabcolsep) * \real{0.1250}}
  >{\raggedleft\arraybackslash}p{(\linewidth - 14\tabcolsep) * \real{0.1250}}
  >{\raggedleft\arraybackslash}p{(\linewidth - 14\tabcolsep) * \real{0.1250}}
  >{\raggedleft\arraybackslash}p{(\linewidth - 14\tabcolsep) * \real{0.1250}}
  >{\raggedleft\arraybackslash}p{(\linewidth - 14\tabcolsep) * \real{0.1250}}
  >{\raggedleft\arraybackslash}p{(\linewidth - 14\tabcolsep) * \real{0.1250}}@{}}
\toprule\noalign{}
\begin{minipage}[b]{\linewidth}\raggedright
Period
\end{minipage} & \begin{minipage}[b]{\linewidth}\raggedleft
Total GDP
\end{minipage} & \begin{minipage}[b]{\linewidth}\raggedleft
Industry
\end{minipage} & \begin{minipage}[b]{\linewidth}\raggedleft
Coal
\end{minipage} & \begin{minipage}[b]{\linewidth}\raggedleft
Services
\end{minipage} & \begin{minipage}[b]{\linewidth}\raggedleft
Population
\end{minipage} & \begin{minipage}[b]{\linewidth}\raggedleft
GDP per head
\end{minipage} & \begin{minipage}[b]{\linewidth}\raggedleft
Agriculture
\end{minipage} \\
\midrule\noalign{}
\endhead
\bottomrule\noalign{}
\endlastfoot
1700--1760 & 0.55 & 0.47 & 0.87 & 0.50 & 0.29 & 0.25 & 0.75 \\
1760--1790 & 0.82 & 1.16 & 2.05 & 0.83 & 0.74 & 0.08 & 0.60 \\
1790--1815 & 1.58 & 1.81 & 2.85 & 2.02 & 1.18 & 0.40 & 0.81 \\
1815--1830 & 1.98 & 3.64 & 2.84 & 1.62 & 1.49 & 0.50 & 0.63 \\
1830--1870 & 2.35 & 2.85 & 3.57 & 2.62 & 1.16 & 1.19 & 0.78 \\
\end{longtable}
}

\emph{Source: own calculations from the Bank of England millennium
dataset (Broadberry et al.~2015). Slopes of log-linear trends fitted
within each window.}

Table 2 formalises the comparison. Fixing the break at 1761 and
estimating over 1700--1830, the trend slope rises by between 0.83
(population) and 3.24 (iron) percentage points a year for every
canal-served series, all with p-values below 0.01 under Newey--West
errors. The slope of income per head changes by 0.02 points and that of
agriculture by −0.04, neither distinguishable from zero. The fixed date
is a pre-test, so the table also reports the break the data choose. The
single best break falls between 1777 (population) and 1792 (total
output) for every canal-served series, with sup-F statistics far above
the Andrews critical value, and at 1818 for income per head.
Agriculture's best break is not significant: there is none. Allowing two
breaks places the first in 1774--1775 for total output, industry and
services and the second in 1818--1823 for output and industry; coal's
first break falls earlier, in 1741, and its second in 1798, which we
return to below. The data pick out the two regimes without being told
where to look.

\textbf{Table 2: Change in trend slope at 1761 and data-chosen break
dates}

{\def\LTcaptype{none} % do not increment counter
\begin{longtable}[]{@{}
  >{\raggedright\arraybackslash}p{(\linewidth - 10\tabcolsep) * \real{0.1667}}
  >{\raggedleft\arraybackslash}p{(\linewidth - 10\tabcolsep) * \real{0.1667}}
  >{\raggedleft\arraybackslash}p{(\linewidth - 10\tabcolsep) * \real{0.1667}}
  >{\raggedleft\arraybackslash}p{(\linewidth - 10\tabcolsep) * \real{0.1667}}
  >{\raggedleft\arraybackslash}p{(\linewidth - 10\tabcolsep) * \real{0.1667}}
  >{\raggedright\arraybackslash}p{(\linewidth - 10\tabcolsep) * \real{0.1667}}@{}}
\toprule\noalign{}
\begin{minipage}[b]{\linewidth}\raggedright
Series
\end{minipage} & \begin{minipage}[b]{\linewidth}\raggedleft
Trend before 1761, \% per year
\end{minipage} & \begin{minipage}[b]{\linewidth}\raggedleft
Change after 1761, points per year
\end{minipage} & \begin{minipage}[b]{\linewidth}\raggedleft
p
\end{minipage} & \begin{minipage}[b]{\linewidth}\raggedleft
Best single break (sup-F)
\end{minipage} & \begin{minipage}[b]{\linewidth}\raggedright
Best two breaks
\end{minipage} \\
\midrule\noalign{}
\endhead
\bottomrule\noalign{}
\endlastfoot
Total GDP & 0.52 & +0.85 & \textless0.01 & 1792 (778) & 1775, 1818 \\
Industry & 0.39 & +1.42 & \textless0.01 & 1789 (909) & 1774, 1823 \\
Coal & 0.77 & +1.62 & \textless0.01 & 1784 (667) & 1741, 1798 \\
Iron & 0.26 & +3.24 & \textless0.01 & 1786 (1,294) & 1786, 1815 \\
Textiles & 0.69 & +1.51 & \textless0.01 & 1779 (240) & 1728, 1817 \\
Services & 0.44 & +1.07 & \textless0.01 & 1786 (927) & 1775, 1844 \\
Population & 0.23 & +0.83 & \textless0.01 & 1777 (3,052) & 1730, 1783 \\
\emph{GDP per head} & 0.29 & +0.02 & 0.82 & 1818 (270) & 1720, 1818 \\
\emph{Agriculture} & 0.82 & −0.04 & 0.84 & none (4.2, not significant) &
1720, 1757 \\
\end{longtable}
}

\emph{Slope-change regressions on 1700--1830 with Newey--West (10 lags)
errors. Break searches on 1700--1870: Andrews sup-F with 15 per cent
trimming (5 per cent critical value approximately 11.8 for two
parameters) and a two-break grid search with a minimum segment of 20
years, which cannot place a break after 1850.}

The arithmetic of the first regime is simple. Between 1760 and 1815 the
growth of total output and the growth of population each rose by about a
percentage point a year: the economy grew faster and fed more people at
the same income, for half a century, while coal output per head doubled,
from an index of 100 in 1700 to 202 in 1790 and 246 in 1800. Coal's
early break in 1741 is a reminder that the coal trade did not begin with
canals (Hatcher 1993; Flinn 1984); what the canal era added was inland
coal, and the second break in 1798 dates it.

\subsection{4.2 The canal network as a
dose}\label{the-canal-network-as-a-dose-1}

Figure 2 shows the canal series. Panel (a) gives miles opened per
decade: 117 in the 1760s, 353 in the 1770s, a lull of 83 in the 1780s,
then 548 in the 1790s, 479 in the 1800s and 344 in the 1810s. Panel (b)
gives the cumulative stock, from 218 miles in 1760 to 772 in 1790, 1,487
in 1800 and 2,320 in 1830. Panel (c) gives coal output per head. The
correspondence between the two waves and the two steepenings of coal per
head is visible; the regressions ask whether it survives detrending.

Figure 2: The canal network and coal. (a) Canal miles opened per decade,
155 canals by completion year. (b) Cumulative canal miles since 1700.
(c) Coal output per head, 1700 = 100, with the year steam overtook water
and wind as a source of stationary power. Sources: canal reference table
described in the data section; Bank of England millennium dataset
(Broadberry et al.~2015); horsepower benchmarks from Kanefsky as
tabulated in Crafts (Kanefsky 1979; Crafts 2004).

Table 3 reports the dose-response regressions over 1700--1830. In the
simplest specification, log output on a linear trend and the canal
stock, a thousand miles of canal is associated with 47 per cent more
coal, 40 per cent more industrial output, 99 per cent more iron, 32 per
cent more services and 24 per cent more population, and with no change
in income per head or agriculture. The sceptical specifications thin
this out. With a quadratic trend and a war dummy, coal, iron and
services retain large and significant coefficients; population is
marginal; industry and total output are indistinguishable from the
trend. In first differences with ten lags of new mileage, coal, industry
and population respond, total output and income per head do not, and
agriculture's response is negative. In a horse race against the print
frequency of ``steam'', coal's canal coefficient is 41 per cent against
an insignificant 13 for steam. A horse race against installed steam
horsepower is not informative: over 1760--1830 the interpolated
horsepower series and the canal stock correlate at 0.98, so partial
coefficients apportion a common trend. The ordering evidence must come
from timing.

\textbf{Table 3: Canal stock and British output, 1700--1830 (per 1,000
canal miles)}

{\def\LTcaptype{none} % do not increment counter
\begin{longtable}[]{@{}
  >{\raggedright\arraybackslash}p{(\linewidth - 10\tabcolsep) * \real{0.1667}}
  >{\raggedleft\arraybackslash}p{(\linewidth - 10\tabcolsep) * \real{0.1667}}
  >{\raggedleft\arraybackslash}p{(\linewidth - 10\tabcolsep) * \real{0.1667}}
  >{\raggedleft\arraybackslash}p{(\linewidth - 10\tabcolsep) * \real{0.1667}}
  >{\raggedleft\arraybackslash}p{(\linewidth - 10\tabcolsep) * \real{0.1667}}
  >{\raggedleft\arraybackslash}p{(\linewidth - 10\tabcolsep) * \real{0.1667}}@{}}
\toprule\noalign{}
\begin{minipage}[b]{\linewidth}\raggedright
Outcome (log)
\end{minipage} & \begin{minipage}[b]{\linewidth}\raggedleft
Linear trend
\end{minipage} & \begin{minipage}[b]{\linewidth}\raggedleft
Quadratic trend + war
\end{minipage} & \begin{minipage}[b]{\linewidth}\raggedleft
First differences, 10-year cumulative
\end{minipage} & \begin{minipage}[b]{\linewidth}\raggedleft
Horse race: canal
\end{minipage} & \begin{minipage}[b]{\linewidth}\raggedleft
Horse race: ``steam'' in print
\end{minipage} \\
\midrule\noalign{}
\endhead
\bottomrule\noalign{}
\endlastfoot
Coal & +46.5\% (p\textless0.001) & +29.8\% (p=0.001) & +27.0\%
(p\textless0.001) & +40.8\% (p\textless0.001) & +12.6\% (p=0.08) \\
Industry & +40.1\% (p\textless0.001) & −0.9\% (p=0.89) & +14.3\%
(p\textless0.001) & +28.8\% (p\textless0.001) & +30.4\% (p=0.003) \\
Iron & +99.3\% (p\textless0.001) & +61.2\% (p=0.013) & --- & --- &
--- \\
Services & +32.3\% (p\textless0.001) & +12.2\% (p=0.002) & --- & --- &
--- \\
Population & +24.0\% (p\textless0.001) & +4.1\% (p=0.07) & +16.6\%
(p\textless0.001) & +18.6\% (p\textless0.001) & +15.0\%
(p\textless0.001) \\
Total GDP & +24.7\% (p\textless0.001) & +2.2\% (p=0.52) & +4.2\%
(p=0.25) & +18.8\% (p\textless0.001) & +19.0\% (p\textless0.001) \\
\emph{GDP per head} & +0.8\% (p=0.69) & −1.9\% (p=0.61) & −12.4\%
(p=0.33) & +0.2\% (p=0.94) & +4.0\% (p=0.08) \\
\emph{Agriculture} & −0.5\% (p=0.90) & −5.1\% (p=0.56) & −18.8\%
(p\textless0.001) & −1.5\% (p=0.72) & +7.3\% (p=0.07) \\
\end{longtable}
}

\emph{Newey--West errors with 10 lags. The first-difference column
reports the sum of the coefficients on lags 0--10 of new mileage, with
the p-value of the joint F-test that all eleven are zero. The steam
regressor is the Google Books frequency of ``steam'', indexed to 1830.
The war dummy covers 1756--63, 1775--83 and 1793--1815. An augmented
Dickey--Fuller test rejects a unit root in the residuals of the
linear-trend regressions at 5 per cent for every outcome except
population (p = 0.24).}

Two features of Table 3 matter. Coal is the robust channel, and it is
the channel the mechanism predicts, since the canals were dug to move
it. And the placebo rows behave: income per head does not respond in any
specification, and agriculture's only significant coefficient is
negative, in first differences, the structural shift away from farming
rather than an effect on it.

The reverse regression is not empty: past growth in coal, total output
and population predicts subsequent canal openings, and past coal growth
predicts new parliamentary authorisations with a joint p-value below
0.001. Canals were built where demand was growing. The predetermined
doses address the contemporaneous part of this. With the completion
stock lagged ten or fifteen years the pattern of Table 3 is unchanged:
coal 48 to 49 per cent per thousand miles, industry 40 to 43, population
24 to 25, income per head nothing. With the authorisation count lagged
ten years, each ten canals authorised is followed by 12 per cent more
coal, 10 per cent more industrial output and 6 per cent more population,
again with no response in income per head; adding the outcome's own
growth over the preceding decade leaves these at 11, 9 and 5 per cent,
all significant at 1 per cent. Under a quadratic trend the authorisation
dose falls to 4 per cent for coal, significant only at the 7 per cent
level. In local projections it predicts coal at every horizon and does
not predict population, income per head or agriculture. The dose is not
exogenous; what we claim is that its timing and incidence, under
contemporaneous and predetermined measures alike, are those of a
precondition.

\subsection{4.3 The precondition chain}\label{the-precondition-chain}

Figure 3 reports the local projections that test the sequence directly.
Each panel shows the cumulative response of the outcome, in log points,
to a unit of the regressor, at horizons of one to twenty years, with 95
per cent Newey--West bands whose bandwidth grows with the horizon. The
responses are predictive, not causal.

Figure 3: Local projections along the precondition chain. Cumulative log
response at horizons 1--20 years, per 1,000 canal miles or per log point
of the regressor, controlling for the outcome's level and a trend.
Samples 1700--1830 where the canal stock is the regressor and 1760--1830
where steam horsepower enters, including the steam-to-income panel.
Shaded: 95 per cent bands with Newey--West bandwidth equal to the larger
of ten years and the horizon.

The top-left panel is the first link: a thousand miles of canal raises
coal output by 29 per cent after five years, 38 per cent after ten and
44 per cent after fifteen, with the band excluding zero at every horizon
shown. Controlling for contemporaneous steam horsepower, estimated from
1760, the canal effect on coal remains at 21 per cent over five to ten
years and is gone at fifteen to twenty. The top-middle panel, canal
stock to steam horsepower, is positive and precisely estimated at every
horizon, but because the horsepower series is interpolated between four
benchmarks, the panel establishes that steam capacity rose after the
network did, not how fast. The top-right panel, coal to steam, is
positive at every horizon but significant only at five years. On the
interpolated series that is as much as can be asked.

The bottom row is the test that distinguishes a precondition from a
cause. Canal stock has no effect on income per head at any horizon
(bottom-left); the point estimates are negative and the bands include
zero throughout. Steam horsepower has no effect on income per head when
the sample stops at 1830 (bottom-middle). On the 1760--1870 sample it
does, at every horizon and with p-values below 0.002; estimated on
1830--1870 alone the response is again positive at every horizon with
p-values below 0.001, and an interaction of horsepower with a post-1830
indicator is positive and significant. The elasticities, between 0.2 and
0.9 depending on sample and horizon, should not be read as magnitudes;
the sign and the timing are the result. Steam's per-capita dividend is a
post-1830 phenomenon. Canals raise coal and population (the population
response is 3 to 7 per cent per thousand miles over five to twenty
years, significant beyond five years at the family-wise threshold) and
leave income per head where it was. The bottom-right panel, agriculture,
is negative at intermediate horizons, again the composition effect, and
returns to zero.

Because the horsepower series is interpolated, we re-estimate the steam
links with the print frequency of ``steam engine'', which correlates at
0.90 with interpolated horsepower but varies year to year. The ordering
survives, more weakly: the canal stock predicts ``steam engine'' at ten
years (p = 0.06) and coal output predicts it at twenty (p = 0.003). The
print series predicts income per head negatively at fifteen years, a
warning about the series rather than about steam: it records discussion
of engines, not the power they supplied. The last link of the chain
therefore rests on the horsepower benchmarks, the 1818 break in income
per head and the 1833 crossover in installed power.

The chain, read from Figure 3 and Table 4, is: canals raise coal within
a decade; steam capacity rises after the network and after coal, on a
scale that only the benchmarks can date; steam raises income per head,
but only once it is the majority power source. Water first, coal second,
steam third, income last.

\subsection{4.4 How much power was
steam?}\label{how-much-power-was-steam}

The precondition thesis requires that the first regime run on water
rather than steam. Table 4 uses the Kanefsky benchmarks to check. In
1760 steam supplied about 6 per cent of Britain's installed stationary
power; in 1800, when the first regime was thirty years old, 21 per cent;
by 1820 about a third; in 1830 parity with water, at 47 per cent. Steam
overtook water and wind combined in 1833, fifteen years after the
per-capita break. Coal output per unit of installed steam horsepower
fell from 100 to 37 between 1760 and 1800: coal grew far faster than the
engines that could burn it, and most of it went to hearths, forges,
kilns and salt pans, moved by water.

\textbf{Table 4: Installed stationary power in Britain, thousands of
horsepower}

{\def\LTcaptype{none} % do not increment counter
\begin{longtable}[]{@{}
  >{\raggedright\arraybackslash}p{(\linewidth - 10\tabcolsep) * \real{0.1667}}
  >{\raggedleft\arraybackslash}p{(\linewidth - 10\tabcolsep) * \real{0.1667}}
  >{\raggedleft\arraybackslash}p{(\linewidth - 10\tabcolsep) * \real{0.1667}}
  >{\raggedleft\arraybackslash}p{(\linewidth - 10\tabcolsep) * \real{0.1667}}
  >{\raggedleft\arraybackslash}p{(\linewidth - 10\tabcolsep) * \real{0.1667}}
  >{\raggedleft\arraybackslash}p{(\linewidth - 10\tabcolsep) * \real{0.1667}}@{}}
\toprule\noalign{}
\begin{minipage}[b]{\linewidth}\raggedright
Year
\end{minipage} & \begin{minipage}[b]{\linewidth}\raggedleft
Steam
\end{minipage} & \begin{minipage}[b]{\linewidth}\raggedleft
Water
\end{minipage} & \begin{minipage}[b]{\linewidth}\raggedleft
Wind
\end{minipage} & \begin{minipage}[b]{\linewidth}\raggedleft
Steam share
\end{minipage} & \begin{minipage}[b]{\linewidth}\raggedleft
Coal output per steam hp (1760 = 100)
\end{minipage} \\
\midrule\noalign{}
\endhead
\bottomrule\noalign{}
\endlastfoot
1760 & 5 & 70 & 10 & 6\% & 100 \\
1800 & 35 & 120 & 15 & 21\% & 37 \\
1830 & 160 & 160 & 20 & 47\% & 19 \\
1870 & 2,060 & 230 & 10 & 90\% & 6 \\
\end{longtable}
}

\emph{Source: Kanefsky's estimates as tabulated in Crafts (Kanefsky
1979; Crafts 2004); coal output from the Bank of England millennium
dataset (Broadberry et al.~2015). Shares from log-linear interpolation
between benchmarks, wind included.}

\subsection{4.5 The semantic sequence}\label{the-semantic-sequence}

If coal moved by water before it burned in engines, the language of the
period should say so. Table 5 records, for each term or group, the first
year in which its smoothed frequency in the British corpus reached 10,
25 and 50 per cent of its 1850 level. ``Canal'' reaches a quarter of its
mid-century frequency in 1763, the year after the Bridgewater opening;
``coal barge'' in 1781; ``coal wharf'' in 1800; ``canal boat'', a later
coinage, in 1823. ``Steam engine'' reaches the same threshold in 1808
and ``steam power'' in 1826. The coal-by-water group as a whole crosses
25 per cent in 1800, the coal-by-steam group in 1819. ``Steam engine''
overtakes ``fire engine'', the older name for the same machine, in 1800,
which dates the terminological consolidation of steam to the decade
after the canal mania. Figure 4 plots the three indices.

\textbf{Table 5: Year in which print frequency first reaches a share of
its 1850 level (British English corpus, five-year mean)}

{\def\LTcaptype{none} % do not increment counter
\begin{longtable}[]{@{}lrrr@{}}
\toprule\noalign{}
Term or group & 10\% & 25\% & 50\% \\
\midrule\noalign{}
\endhead
\bottomrule\noalign{}
\endlastfoot
``canal'' & 1743 & 1763 & 1809 \\
``coal barge'' & 1753 & 1781 & 1783 \\
``coal wharf'' & 1774 & 1800 & 1815 \\
``canal boat'' & 1812 & 1823 & 1825 \\
Coal-by-water group & 1780 & 1800 & 1817 \\
``steam engine'' & 1800 & 1808 & 1821 \\
``steam power'' & 1821 & 1826 & 1831 \\
Coal-by-steam group & 1804 & 1819 & 1825 \\
``coal field'' & 1788 & 1820 & 1831 \\
\end{longtable}
}

\emph{Source: Google Books Ngram corpus, British English 2019 edition,
three-year API smoothing and a five-year centred mean. Groups are
equal-weighted averages of member terms each indexed to 1850; the
coal-by-water group also contains ``coal boat'', which is too rare
before 1800 to threshold on its own.}

Figure 4: In print, coal travels by water before it is burned in
engines. Five-year moving averages indexed to 1850 = 100. Coal-by-water:
equal-weighted ``coal barge'', ``coal wharf'', ``coal boat'', ``canal
boat''. Coal-by-steam: ``steam engine'', ``steam power''. Source: Google
Books Ngram corpus, British English 2019 edition.

The ordering does not depend on the smoothing or the reference year.
Across three-, five- and nine-year windows and with 1830 or 1850 as
reference, ``coal barge'' crosses a quarter of its reference level in
1780--1781, ``canal'' in 1763--1766, ``steam engine'' in 1807--1809 and
``steam power'' in 1822--1826; the coal-by-water group crosses between
1785 and 1801 and the coal-by-steam group between 1814 and 1819 in every
combination. ``Coal wharf'' is the one sensitive term, crossing in 1776
or 1800 depending on the reference, and in both cases before steam. Read
as decades, the order is stable.

The corpus also tells us what kind of thing the infrastructure
vocabulary measures. The frequency of ``canal'' correlates at 0.91 with
the cumulative mileage of canals in existence over 1740--1850 and at
−0.04 with the mileage opened in the surrounding decade; linearly
detrended, the first correlation falls to 0.23, and in first differences
to 0.16. Print records the level of the network, which is written about
every year it exists, not the rate of building. We use it only in the
first sense.

\subsection{4.6 The comparative cases: water without coal, coal without
water}\label{the-comparative-cases-water-without-coal-coal-without-water}

Tvedt's thesis is comparative, and the British result gives it a form
the comparative cases can be checked against. If engineered water was a
precondition for a coal economy and coal the fuel of the second regime,
then an economy with waterways and no coal should show at most the first
regime; an economy with coal and no engineered water should show neither
until the water arrives; and an economy that acquires both late should
run the sequence compressed. Annual series joining output, canal
building and installed power are not available elsewhere, so the check
is at the Maddison benchmark years, where population is measured rather
than interpolated, read alongside what the national literatures record
about fuel and water. Table 6 gives growth between 1700 and 1820, before
steam supplied a third of British power, and between 1820 and 1870.

\textbf{Table 6: Growth between benchmark years, per cent}

{\def\LTcaptype{none} % do not increment counter
\begin{longtable}[]{@{}
  >{\raggedright\arraybackslash}p{(\linewidth - 8\tabcolsep) * \real{0.2000}}
  >{\raggedleft\arraybackslash}p{(\linewidth - 8\tabcolsep) * \real{0.2000}}
  >{\raggedleft\arraybackslash}p{(\linewidth - 8\tabcolsep) * \real{0.2000}}
  >{\raggedleft\arraybackslash}p{(\linewidth - 8\tabcolsep) * \real{0.2000}}
  >{\raggedleft\arraybackslash}p{(\linewidth - 8\tabcolsep) * \real{0.2000}}@{}}
\toprule\noalign{}
\begin{minipage}[b]{\linewidth}\raggedright
\end{minipage} & \begin{minipage}[b]{\linewidth}\raggedleft
GDP per head 1700--1820
\end{minipage} & \begin{minipage}[b]{\linewidth}\raggedleft
Population 1700--1820
\end{minipage} & \begin{minipage}[b]{\linewidth}\raggedleft
Total GDP 1700--1820
\end{minipage} & \begin{minipage}[b]{\linewidth}\raggedleft
GDP per head 1820--1870
\end{minipage} \\
\midrule\noalign{}
\endhead
\bottomrule\noalign{}
\endlastfoot
Britain (UK, Maddison) & +37 & +148 & +240 & +76 \\
Great Britain (Broadberry) & +37 & +125 & +209 & +76 \\
Germany & +35 & +66 & +124 & +44 \\
Belgium & +8 & +72 & +85 & +82 \\
Spain & +22 & +39 & +69 & +16 \\
France & +3 & +46 & +50 & +65 \\
Sweden & −29 & +104 & +44 & +52 \\
Netherlands & −11 & +23 & +9 & +47 \\
China & −43 & +176 & +58 & +7 \\
\end{longtable}
}

\emph{Sources: Maddison Project Database 2023 (Bolt and van Zanden
2025); Great Britain row from the Bank of England millennium dataset
(Broadberry et al.~2015). The 1700 figure for China is contested
(Broadberry, Guan, and Li 2018).}

On either measure British output roughly tripled between 1700 and 1820,
and Britain was the only economy in the panel to more than double its
population while raising income per head; Germany doubled its output and
France grew by half.

The Netherlands is the case of water without coal. By the 1660s it had
built more than six hundred kilometres of trekvaarten, towpath canals
with scheduled barge services, on top of the densest river and drainage
network in Europe (de Vries 1978). It ran on peat, which gave it in the
seventeenth century an energy supply per head no other European economy
matched, and it imported the coal it burned from Newcastle and Liège (de
Zeeuw 1978; Kander, Malanima, and Warde 2013). What it lacked, besides
coal, was gradient: a country at sea level has water to float on and
little to fall, and Dutch industry was powered by wind. Its total output
grew 9 per cent between 1700 and 1820 and its income per head fell, for
reasons the literature locates in trade, finance, wages and war rather
than fuel (de Vries and van der Woude 1997). The case shows that water
transport was not sufficient; it does not, on its own, show that coal
was what was missing. What it shows is that the densest engineered
waterways in Europe, with no cheap domestic fuel to carry and no falling
water to drive mills, did not compound.

Belgium is the case of coal that waited for water, and it is the nearest
thing in the panel to an experiment on the sequence. The Hainaut and
Liège coalfields were worked throughout the eighteenth century, but the
Mons--Condé canal that joined the Borinage to the Scheldt opened in 1818
and the Charleroi--Brussels canal in 1832 (Mokyr 1976). Belgian income
per head grew 8 per cent between 1700 and 1820, less than Spain's, and
82 per cent between 1820 and 1870, more than Britain's, on the way to
the densest railway network on the continent. Belgium ran Britain's
sequence in a third of the time, and it began after the canals reached
the pits. The second regime could arrive so quickly because steam, by
1820, was a technology that could be imported; the first regime could
not be, because it was built into the ground.

China is the case of coal and water that did not meet. The Grand Canal,
some 1,800 kilometres from Hangzhou to Beijing, was the longest
artificial waterway in the world and a state conduit for grain (Elvin
1973). China's coal lay in the north-west, in Shanxi and Shaanxi, a
thousand kilometres from the Jiangnan core where the population, the
textile industry and the demand for fuel were (Pomeranz 2000). The
rivers that might have joined them were of the wrong kind: the Yellow
River carried silt rather than barges and changed course, and the
monsoon gave rivers and canals alike a wet season and a dry one (Tvedt
2010). The benchmarks record a 176 per cent rise in population and a 43
per cent fall in income per head between 1700 and 1820, though the 1700
level is contested (Broadberry, Guan, and Li 2018), and 7 per cent
per-capita growth in the half-century after. This is what a Malthusian
regime looks like, and it is the regime Britain left in aggregate in the
1770s and per head in 1818: an economy that had exhausted its organic
frontier and could not reach its mineral one.

India we treat in a sentence, on Parthasarathi's evidence. Its coal lay
in Bengal, at Raniganj, while its export industry, cotton, lay in
Gujarat and on the Coromandel coast, and the rivers between ran with the
monsoon (Parthasarathi 2011; Tvedt 2010). The endowment was China's
problem in a different geography.

None of these cases is a test in the sense that the British series are.
They are four readings of one rule, and the rule is what the British
result supplies: the economies with engineered water and no coal, or
coal and no engineered water, did not leave the organic regime in this
period; the two that had both, Britain from the 1770s and Belgium from
the 1820s, did, in the order the precondition predicts. A comparative
paper with regional series for the Low Countries and the Yangzi delta
could turn the readings into tests.

\subsection{4.7 Why we stop at benchmark
years}\label{why-we-stop-at-benchmark-years}

The natural cross-country test of an eighteenth-century British take-off
is a difference-in-differences on Maddison GDP per head, treating
Britain from 1761 against continental controls. In logs, with France and
the Netherlands as controls and year and country effects, it returns
treatment coefficients of 0.08 to 0.31 with p-values below 0.001 for
every control group and window we tried. Taken at face value they say
that Britain pulled away from the continent from 1761. They should not
be. Britain's gap to the controls is significantly negative in the 1700s
and 1710s, so there is no parallel pre-trend to break. A sup-F search on
the log gap to the Netherlands and France places the single break in
1807, and Figure 5 shows what happened then: Britain's income per head,
indexed to 1790, stood at 109 in 1805 and 114 in 1815; the Netherlands'
at 100 in 1805, 63 in 1808 and 72 in 1815. Underneath both, Britain's
own income per head grew at 0.08 per cent a year between 1760 and 1790.
There was no canal-era per-capita acceleration for any control group to
reveal.

Table 7 gives the drawdowns. Between the late 1780s and the Napoleonic
trough the Netherlands lost 44 per cent of its income per head, Portugal
48, Sweden 27, France 22 and Spain 13. Britain lost 1.4 per cent. The
1761 ``treatment effect'' measured against a continental control group
is, to a first approximation, the difference between being blockaded and
doing the blockading. That is why the comparative evidence above is read
at benchmark years and for aggregate output: over 1790--1815 the
per-capita series record the war, not the rivers.

\textbf{Table 7: GDP per head, peak 1785--95 to trough 1795--1815
(Maddison Project Database 2023, 2011 international dollars)}

{\def\LTcaptype{none} % do not increment counter
\begin{longtable}[]{@{}lrrrr@{}}
\toprule\noalign{}
Country & Peak & Trough & Drawdown & 1815 relative to 1790 \\
\midrule\noalign{}
\endhead
\bottomrule\noalign{}
\endlastfoot
Britain (UK) & 3,207 (1795) & 3,161 (1798) & −1.4\% & +13.6\% \\
Netherlands & 4,666 (1794) & 2,632 (1808) & −43.6\% & −28.3\% \\
Portugal & 2,063 (1785) & 1,072 (1811) & −48.0\% & −24.1\% \\
Sweden & 1,661 (1791) & 1,221 (1809) & −26.5\% & −11.9\% \\
France & 2,016 (1788) & 1,580 (1801) & −21.7\% & +1.3\% \\
Spain & 1,454 (1790) & 1,265 (1811) & −13.0\% & +3.2\% \\
Germany & 1,820 (1792) & 1,725 (1805) & −5.2\% & +8.1\% \\
\end{longtable}
}

\emph{Source: Maddison Project Database 2023 (Bolt and van Zanden 2025).
The series labelled Britain is the United Kingdom, including Ireland.}

Figure 5: The 1807 ``break'' in the cross-country design is the Dutch
collapse. GDP per head, 1790 = 100, for Britain, France and the
Netherlands, with the Revolutionary and Napoleonic war years shaded and
the estimated break in the Britain--controls gap marked. Source:
Maddison Project Database 2023 (Bolt and van Zanden 2025).

\begin{center}\rule{0.5\linewidth}{0.5pt}\end{center}

\section{5. Discussion}\label{discussion}

\subsection{5.1 Floated before fired}\label{floated-before-fired}

The results describe an economy that changed gear twice. The first
change, in the 1770s and 1780s, was in the quantity of things: coal dug,
iron smelted, goods made, people fed. The second, after 1818, was in the
quantity per person. The first coincided with the building of the canal
network and ran on water power and horse-drawn barges; the second
coincided with steam becoming the majority power source. Coal is the
hinge between them. Its output per head doubled during the first regime,
when steam supplied between a twentieth and a fifth of installed power,
and it is the one outcome whose response to the canal stock survives
every specification, including the most demanding, where it retains
significance only at the 7 per cent level.

This is Tvedt's sequence in the annual record. Britain's water system
did not compete with coal; it delivered coal. The Bridgewater Canal
existed to bring the Worsley pits to Manchester; the Grand Cross, begun
in 1766 and completed when the Oxford Canal reached the Thames in 1790,
joined the Trent, Mersey, Severn and Thames; the mania canals of the
1790s ran, disproportionately, from coalfields to towns. The sea-coal
trade to London carried more tonnage throughout the century; what canals
changed was where coal could go inland, and at what price (Pollard
1980). By 1800 coal was cheap far from the pithead, and the population
had grown by half over the same forty years. We do not claim that canals
caused that growth, which the demographic literature attributes to
earlier marriage (Wrigley 1988); we claim that the economy fed it at
constant income, and that cheap inland coal and freight were part of
how. The steam engine then found, ready-made, a market, a fuel supply, a
factory system organised around water wheels, and in the canal company a
legal template for raising capital by Act of Parliament (Ward 1974). It
had to improve on a prime mover that was, in 1800, three times its size.

\subsection{5.2 What the first regime was, and was
not}\label{what-the-first-regime-was-and-was-not}

The first regime was not a rise in living standards. Income per head did
not move, and the local projections show that canals did not move it. It
was a rise in carrying capacity: the same income for half again as many
people, sustained for two generations without the Malthusian check that
had ended every earlier expansion. In Wrigley's terms this is not an
escape from the organic economy, which by his definition Britain had
begun to leave when it turned to coal for heat in the sixteenth century
(Wrigley 1988; Warde 2007); it is the advanced organic economy working
at full stretch, with mineral heat, organic power and engineered water.
The escape from the ceiling in aggregate came before the escape per
head, and it came before steam.

This also disposes of a false choice. Malm argues that the transition
from water to steam power in the 1830s and 1840s was driven by the
spatial and disciplinary needs of capital rather than by the exhaustion
of water (Malm 2016). Our results are consistent with that and add a
prior stage to it: before steam could be chosen over water for its
portability, water had to have created the markets and the fuel supply
that made a portable prime mover worth having. Allen's
induced-innovation story likewise needs coal to be cheap where engines
were built (Allen 2009); the canals are a large part of why it was.

\subsection{5.3 The two clocks
reconciled}\label{the-two-clocks-reconciled}

The cliometric revision of the 1980s established that income per head
grew slowly until 1830 and has been read as denying an
eighteenth-century take-off (Crafts 1985; Crafts and Harley 1992). Our
results accept the revision and deny the reading. There was a take-off
in the eighteenth century; it was an aggregate one, and it was absorbed
by population. The canal historians and the growth accountants have been
measuring different variables and talking past each other. Crafts's
finding that steam's contribution to productivity growth is negligible
before 1830 and large after 1850 is the second regime seen from the
supply side (Crafts 2004). Our finding that steam horsepower raises
income per head only after 1830 is the same regime seen from the demand
side.

One rival explanation of the 1818 break deserves to be named, because we
have used it against the cross-country design: peace. Demobilisation,
deflation and reopened trade could produce a per-capita break in
1816--1818 without steam. Absorbing a war dummy moves the estimated
break from 1818 to 1816 and does not remove it, and the per-capita trend
over 1816--1830, at 0.7 per cent a year, is well below the 1.2 per cent
of 1830--1870. The break is real; its timing is over-determined; and the
acceleration that follows it belongs to the steam decades.

\subsection{5.4 Tvedt's question, answered as far as Britain can answer
it}\label{tvedts-question-answered-as-far-as-britain-can-answer-it}

The comparative readings in section 4.6 let us say what the British test
does for the global question. Tvedt's thesis, as we have restated it, is
a rule about matching. Engineered water was necessary for a coal economy
because, before railways, it was the only way to make coal cheap far
from the pit; coal was necessary for the escape from the organic ceiling
because water and wind could not be scaled. The British series establish
the sequence in the one case we can watch year by year. In the benchmark
cases no economy skipped a step. The Netherlands had the water and
stayed organic. China had the coal and stayed organic. Belgium got the
water in 1818 and left within a generation. Pomeranz's question, why
Britain and not the Yangzi delta, and Tvedt's, why Britain and not the
Netherlands, have on this reading one answer with the terms reversed:
the delta lacked water that could reach its coal, and the Netherlands
lacked coal its water could reach (Pomeranz 2000; Tvedt 2010).

What the rule does not do is explain why Britain matched the two when it
did. That Britain's rivers could be engineered at low cost is Tvedt's
point; that they were engineered in the 1760s rather than the 1660s is a
question about landowners, Parliament and the price of coal in
Manchester, on which national series are silent (Ward 1974). A rule
about necessary conditions leaves the timing of the sufficient ones
open, and four cases cannot exclude other conditions, institutions or
Atlantic trade among them, that would also fit them.

\subsection{5.5 A caution for historical
difference-in-differences}\label{a-caution-for-historical-difference-in-differences}

The cross-country result is methodological. In the long eighteenth
century the control group went to war on its own territory and the
treatment group did not, so a comparison spanning 1790--1815 attributes
the difference to whatever the British treatment happens to be. The
remedies are the ones applied above: locate the break before
interpreting the coefficient, examine the controls' own series, and
prefer within-country evidence where the treatment is national. The
estimator is not at fault; the history is (Roth et al.~2023).

\subsection{5.6 What the text can and cannot
measure}\label{what-the-text-can-and-cannot-measure}

The most useful text result is the negative one: the corpus records the
level of the built environment, not its growth, so vocabulary is a fair
proxy for stock and a poor one for investment. The bigram sequencing is
on firmer ground because it asks about order rather than magnitude.
``Coal barge'' around 1781 and ``coal wharf'' around 1800 record coal on
the water; ``steam engine'' around 1808 and ``steam power'' around 1826
record coal in the engine. The order in the language is the order in the
economy, give or take a decade, and the technical-publishing bias of the
corpus works in favour of exactly this kind of vocabulary (Pechenick et
al.~2015).

\begin{center}\rule{0.5\linewidth}{0.5pt}\end{center}

\section{6. Limitations}\label{limitations}

The argument rests on timing, sequence and incidence in a single
national time series, and it carries the limitations of that design.

\subsection{6.1 The dose is endogenous}\label{the-dose-is-endogenous}

Canals were built where coal and people were. The reverse regressions
confirm it: past growth in coal, total output or population predicts
subsequent canal openings, and past coal growth predicts new
authorisations. The dose-response coefficients are therefore not causal
effects of an exogenous treatment. The predetermined doses remove
contemporaneous feedback, and adding the outcome's own past growth
leaves the authorisation coefficients intact, but nothing here removes
anticipation: a promoter who expected coal demand to grow would seek an
Act in advance of it. What the design establishes is that the canal
stock precedes coal output by five to twenty years under every measure
of the dose and that income per head does not respond. Turnbull's
regional evidence and Alvarez-Palau and co-authors' market-access
estimates supply the cross-sectional identification this paper lacks
(Turnbull 1987; Alvarez-Palau et al.~2025); a county panel with terrain
as the source of exogenous variation in canal access is the natural next
step.

\subsection{6.2 Trending series and multiple
tests}\label{trending-series-and-multiple-tests}

Every series in the paper trends. The level regressions in Table 3 are
exposed to spurious regression; the augmented Dickey--Fuller test on
their residuals rejects a unit root at 5 per cent for every outcome
except population, so the level evidence for population should be
discounted and the first-difference and local-projection results relied
on. The fixed 1761 break is a pre-test; the sup-F statistics are
reported against the Andrews critical value, and agriculture's is not
significant. With eight outcomes, four specifications and four horizons,
a Bonferroni threshold of about 0.002 is the honest bar: coal, industry,
population beyond five years and the post-1830 steam result clear it;
the five-year population response, the authorisation dose under a
quadratic trend and the coal-to-steam link do not.

\subsection{6.3 The canal series is
provisional}\label{the-canal-series-is-provisional}

The mileage series is built from a published reference table of 155
canals by completion year. It omits river navigations improved before
1700 and minor branches, and it dates staged openings to their
completion. Priestley's authorisation dates corroborate the two waves
independently, but they were recovered by optical character recognition
from an 1831 text and only 152 of 325 entries yielded a parseable year.
It excludes the Scottish canals, some 150 miles by 1830, against output
series for Great Britain. A definitive series would come from the
Cambridge Group's sectional opening dates, which are not publicly
deposited (Satchell 2017). Our magnitudes per thousand miles should be
read with that in mind; the timing is less fragile.

\subsection{6.4 Steam and water horsepower are
interpolated}\label{steam-and-water-horsepower-are-interpolated}

Kanefsky's estimates exist at four dates within our sample. The
interpolated series are log-linear between them, so regressions that use
them are identified from three kinks: the local projections involving
steam establish ordering across benchmarks rather than annual dynamics,
their standard errors understate the uncertainty, and the elasticities
of industry and income per head with respect to steam should not be read
as magnitudes. The same interpolation makes water horsepower collinear
with the canal stock over 1760--1830 (r = 0.98), so the two halves of
Tvedt's water thesis, transport and power, cannot be separated with
these data; our dose is transport only. The print frequency of ``steam
engine'' confirms the canal-to-steam ordering weakly and cannot check
the steam-to-income link; engine counts by decade would tighten both
(Kanefsky and Robey 1980). That steam raises income per head only after
1830 does not depend on the interpolation: it follows from the 1818
break, the 1833 crossover in installed power, and estimation on
1830--1870 alone.

\subsection{6.5 Agriculture is an imperfect
placebo}\label{agriculture-is-an-imperfect-placebo}

Agriculture's response to the canal stock is negative and significant at
ten to fifteen years in the local projections and in first differences.
We interpret this as compositional: labour and capital moved out of
farming as canal-served sectors grew. It is not, strictly, a null
result, and a reader who prefers a cleaner placebo can substitute income
per head, which is null in every specification.

\subsection{6.6 Text frequencies measure print, not
speech}\label{text-frequencies-measure-print-not-speech}

The Google Books corpus over-represents technical and legal publishing,
and its British-English sub-corpus is small in the eighteenth century,
so bigram frequencies before 1770 are noisy; the jump of ``coal barge''
from a quarter to half of its reference level in two years is a count
spike, not an event. We therefore use thresholds of a reference level
rather than absolute frequencies, and groups of terms rather than single
terms. The ordering of water terms before steam terms holds for every
smoothing window and reference year we tried; individual crossing years
move by one or two years, and ``coal wharf'' by more, so the years
should be read as decades. ``Fire engine'' also denotes a fire pump,
which is why we use it only to date the terminological consolidation of
``steam engine''.

\subsection{6.7 The comparative cases are readings, not
tests}\label{the-comparative-cases-are-readings-not-tests}

The Dutch, Belgian, Chinese and Indian cases rest on Maddison benchmarks
and on the national literatures, not on series comparable to the British
ones. Annual Dutch output by sector exists from the fourteenth century,
but the trekvaart network was complete by 1665 and supplies no varying
canal dose. Population in the Maddison database is interpolated between
1700 and 1820, so cross-country regressions on total output over that
period are not legitimate and we have not run them; Table 6 is
descriptive, and its Maddison row for Britain is the United Kingdom
including Ireland. The trekvaart mileage, the Belgian canal dates and
the location of Chinese and Indian coal are taken from the works cited,
and the 1700 Chinese benchmark is contested. The war-drawdown result is
a property of the controls and so is robust to their choice, but it
means that cross-country per-capita evidence for a British take-off
before 1815 is, in our reading, unrecoverable from the Maddison series
with any difference-in-differences design.

\begin{center}\rule{0.5\linewidth}{0.5pt}\end{center}

\section{7. Conclusion}\label{conclusion}

Why England, and not the Netherlands, China or India? The answer this
paper can give is narrower than Tvedt's and better measured. Britain
industrialised on two clocks. The first, set in the 1770s, measured how
much the economy produced and how many people it fed; the second, set
around 1818, measured how much each of them had. The first ran on water.
Canals carried the coal inland, water wheels drove the mills, and the
aggregate economy accelerated while steam supplied a twentieth and then
a fifth of its power. The second ran on steam, and it was steam that
finally raised income per head. Coal connects the two: its output per
head doubled in the first regime, when most of it was burned in hearths,
forges and kilns rather than engines, and it is the one variable whose
response to the canal network survives every specification we tried.

This is what a precondition looks like in data. Water infrastructure did
not cause the Industrial Revolution in the sense of raising British
incomes; it did not, and the local projections say so. It made the coal
economy possible by making coal cheap far from the pit, and it helped
feed, at constant income for two generations, the population that would
work in the steam economy. Read as a rule, the British sequence says
that neither engineered water nor coal was sufficient and that the
fossil economy began where the two met. The Netherlands, with the water
and without the coal, stayed in the organic regime; China, with the coal
and without water that could reach it, stayed there too, and India with
it; Belgium, which had the coal and got the water in 1818, ran Britain's
sequence in a third of the time. Four benchmark cases cannot prove a
rule, but they can fail to contradict it, and they do.

Three things follow for how the period is studied. The distinction
between aggregate and per-capita growth should be made explicit in any
account of the eighteenth century, because the two chronologies that
have divided the literature are chronologies of different variables.
Cross-country difference-in-differences spanning the Revolutionary and
Napoleonic wars should be read with the controls' own histories in view,
because the largest economic event of the period happened to the control
group. And the digitised print record, used with care, can date the use
of the built environment even where it cannot measure its growth; in
this case it dates the coal barge a generation before the steam engine.

The fossil economy was floated before it was fired. That is a claim
about order; the order is in the British data, and the comparative
record is consistent with it.

\begin{center}\rule{0.5\linewidth}{0.5pt}\end{center}

\section{8. References}\label{references}

Allen, Robert C. 2009. \emph{The British Industrial Revolution in Global
Perspective}. Cambridge: Cambridge University Press.

Alvarez-Palau, Eduard J., Dan Bogart, Max Satchell, and Leigh
Shaw-Taylor. 2025. ``Transport and Urban Growth in the First Industrial
Revolution.'' \emph{The Economic Journal} 135 (668): 1191--1228.
https://doi.org/10.1093/ej/ueae111.

Andrews, Donald W. K. 1993. ``Tests for Parameter Instability and
Structural Change with Unknown Change Point.'' \emph{Econometrica} 61
(4): 821--856.

Bai, Jushan, and Pierre Perron. 1998. ``Estimating and Testing Linear
Models with Multiple Structural Changes.'' \emph{Econometrica} 66 (1):
47--78.

Bertrand, Marianne, Esther Duflo, and Sendhil Mullainathan. 2004. ``How
Much Should We Trust Differences-in-Differences Estimates?''
\emph{Quarterly Journal of Economics} 119 (1): 249--275.

Bogart, Dan. 2014. ``The Transport Revolution in Industrialising
Britain: A Survey.'' In \emph{The Cambridge Economic History of Modern
Britain, Volume 1: 1700--1870}, edited by Roderick Floud, Jane
Humphries, and Paul Johnson, 368--391. Cambridge: Cambridge University
Press. https://doi.org/10.1017/CHO9781139815017.014.

Bolt, Jutta, and Jan Luiten van Zanden. 2025. ``Maddison-Style Estimates
of the Evolution of the World Economy: A New 2023 Update.''
\emph{Journal of Economic Surveys} 39 (2): 631--671.
https://doi.org/10.1111/joes.12618.

Broadberry, Stephen, Hanhui Guan, and David Daokui Li. 2018. ``China,
Europe, and the Great Divergence: A Study in Historical National
Accounting, 980--1850.'' \emph{Journal of Economic History} 78 (4):
955--1000.

Broadberry, Stephen, Bruce M. S. Campbell, Alexander Klein, Mark
Overton, and Bas van Leeuwen. 2015. \emph{British Economic Growth,
1270--1870}. Cambridge: Cambridge University Press.

Clark, Gregory, and David Jacks. 2007. ``Coal and the Industrial
Revolution, 1700--1869.'' \emph{European Review of Economic History} 11
(1): 39--72.

Crafts, Nicholas F. R. 1985. \emph{British Economic Growth during the
Industrial Revolution}. Oxford: Clarendon Press.

Crafts, Nicholas F. R. 2004. ``Steam as a General Purpose Technology: A
Growth Accounting Perspective.'' \emph{The Economic Journal} 114 (495):
338--351. https://doi.org/10.1111/j.1468-0297.2003.00200.x.

Crafts, Nicholas F. R., and C. Knick Harley. 1992. ``Output Growth and
the British Industrial Revolution: A Restatement of the Crafts--Harley
View.'' \emph{Economic History Review} 45 (4): 703--730.

Crouzet, François. 1964. ``Wars, Blockade, and Economic Change in
Europe, 1792--1815.'' \emph{Journal of Economic History} 24 (4):
567--588.

de Vries, Jan.~1978. \emph{Barges and Capitalism: Passenger
Transportation in the Dutch Economy, 1632--1839}. Utrecht: HES
Publishers.

de Vries, Jan, and Ad van der Woude. 1997. \emph{The First Modern
Economy: Success, Failure, and Perseverance of the Dutch Economy,
1500--1815}. Cambridge: Cambridge University Press.

de Zeeuw, J. W. 1978. ``Peat and the Dutch Golden Age: The Historical
Meaning of Energy-Attainability.'' \emph{A.A.G. Bijdragen} 21: 3--31.

Elvin, Mark. 1973. \emph{The Pattern of the Chinese Past: A Social and
Economic Interpretation}. Stanford: Stanford University Press.

Fernihough, Alan, and Kevin Hjortshøj O'Rourke. 2021. ``Coal and the
European Industrial Revolution.'' \emph{The Economic Journal} 131 (635):
1135--1149.

Flinn, Michael W. 1984. \emph{The History of the British Coal Industry,
Volume 2: 1700--1830, The Industrial Revolution}. Oxford: Clarendon
Press.

Hadfield, Charles. 1984. \emph{British Canals: An Illustrated History}.
7th ed.~Newton Abbot: David and Charles.

Hatcher, John. 1993. \emph{The History of the British Coal Industry,
Volume 1: Before 1700}. Oxford: Clarendon Press.

Jordà, Òscar. 2005. ``Estimation and Inference of Impulse Responses by
Local Projections.'' \emph{American Economic Review} 95 (1): 161--182.

Kander, Astrid, Paolo Malanima, and Paul Warde. 2013. \emph{Power to the
People: Energy in Europe over the Last Five Centuries}. Princeton:
Princeton University Press.

Kanefsky, John W. 1979. ``The Diffusion of Power Technology in British
Industry, 1760--1870.'' PhD thesis, University of Exeter.

Kanefsky, John, and John Robey. 1980. ``Steam Engines in 18th-Century
Britain: A Quantitative Assessment.'' \emph{Technology and Culture} 21
(2): 161--186.

Landes, David S. 1969. \emph{The Unbound Prometheus: Technological
Change and Industrial Development in Western Europe from 1750 to the
Present}. Cambridge: Cambridge University Press.

Maw, Peter. 2013. \emph{Transport and the Industrial City: Manchester
and the Canal Age, 1750--1850}. Manchester: Manchester University Press.

Malm, Andreas. 2016. \emph{Fossil Capital: The Rise of Steam Power and
the Roots of Global Warming}. London: Verso.

Michel, Jean-Baptiste, Yuan Kui Shen, Aviva Presser Aiden, Adrian Veres,
Matthew K. Gray, and Erez Lieberman Aiden. 2011. ``Quantitative Analysis
of Culture Using Millions of Digitized Books.'' \emph{Science} 331
(6014): 176--182.

Mokyr, Joel. 1976. \emph{Industrialization in the Low Countries,
1795--1850}. New Haven: Yale University Press.

Newey, Whitney K., and Kenneth D. West. 1987. ``A Simple, Positive
Semi-Definite, Heteroskedasticity and Autocorrelation Consistent
Covariance Matrix.'' \emph{Econometrica} 55 (3): 703--708.

Parthasarathi, Prasannan. 2011. \emph{Why Europe Grew Rich and Asia Did
Not: Global Economic Divergence, 1600--1850}. Cambridge: Cambridge
University Press.

Pechenick, Eitan Adam, Christopher M. Danforth, and Peter Sheridan
Dodds. 2015. ``Characterizing the Google Books Corpus: Strong Limits to
Inferences of Socio-Cultural and Linguistic Evolution.'' \emph{PLOS ONE}
10 (10): e0137041.

Pollard, Sidney. 1980. ``A New Estimate of British Coal Production,
1750--1850.'' \emph{Economic History Review} 33 (2): 212--235.

Satchell, Max. 2017. ``Navigable Waterways and the Economy of England
and Wales, 1600--1835.'' Cambridge Group for the History of Population
and Social Structure, Online Historical Atlas of Transport.

Pomeranz, Kenneth. 2000. \emph{The Great Divergence: China, Europe, and
the Making of the Modern World Economy}. Princeton: Princeton University
Press.

Priestley, Joseph. 1831. \emph{Historical Account of the Navigable
Rivers, Canals, and Railways, throughout Great Britain}. London:
Longman, Rees, Orme, Brown and Green.

Rambachan, Ashesh, and Jonathan Roth. 2023. ``A More Credible Approach
to Parallel Trends.'' \emph{Review of Economic Studies} 90 (5):
2555--2591.

Roth, Jonathan, Pedro H. C. Sant'Anna, Alyssa Bilinski, and John Poe.
2023. ``What's Trending in Difference-in-Differences? A Synthesis of the
Recent Econometrics Literature.'' \emph{Journal of Econometrics} 235
(2): 2218--2244.

Szostak, Rick. 1991. \emph{The Role of Transportation in the Industrial
Revolution: A Comparison of England and France}. Montreal:
McGill-Queen's University Press.

Thomas, Ryland, and Nicholas Dimsdale. 2017. \emph{A Millennium of UK
Data}. Bank of England OBRA dataset.
https://www.bankofengland.co.uk/statistics/research-datasets.

Turnbull, Gerard. 1987. ``Canals, Coal and Regional Growth during the
Industrial Revolution.'' \emph{Economic History Review} 40 (4):
537--560.

Tvedt, Terje. 2010. ``Why England and Not China and India? Water Systems
and the History of the Industrial Revolution.'' \emph{Journal of Global
History} 5 (1): 29--50.

Ward, J. R. 1974. \emph{The Finance of Canal Building in
Eighteenth-Century England}. Oxford: Oxford University Press.

Warde, Paul. 2007. \emph{Energy Consumption in England and Wales,
1560--2000}. Naples: CNR-ISSM.

Wrigley, E. A. 1988. \emph{Continuity, Chance and Change: The Character
of the Industrial Revolution in England}. Cambridge: Cambridge
University Press.

Wrigley, E. A. 2010. \emph{Energy and the English Industrial
Revolution}. Cambridge: Cambridge University Press.

Wrigley, E. A. 2016. \emph{The Path to Sustained Growth: England's
Transition from an Organic Economy to an Industrial Revolution}.
Cambridge: Cambridge University Press.

\begin{center}\rule{0.5\linewidth}{0.5pt}\end{center}

\section{Data Availability Statement}\label{data-availability-statement}

All code and data required to reproduce the analyses are publicly
available at \url{https://github.com/percw/water_and_society}. British
sectoral output, population and capital stock are from the Bank of
England's \emph{A Millennium of Macroeconomic Data for the UK} (Thomas
and Dimsdale 2017), which reproduces Broadberry et al.~(2015).
Cross-country GDP per head and population are from the Maddison Project
Database 2023 (Bolt and van Zanden 2025). Installed horsepower
benchmarks are from Kanefsky (1979) as reported in Crafts (2004). Canal
completion years and lengths are compiled from published reference
tables and parliamentary authorisation years are parsed from the
digitised text of Priestley (1831); both series are included in the
repository with their construction scripts. Word and phrase frequencies
are from the Google Books Ngram Corpus, British English 2019 edition. A
self-contained replication package is available as a supplementary
archive.

\begin{center}\rule{0.5\linewidth}{0.5pt}\end{center}

\end{document}
